\noindent\textbf{Operaciones y resultados intermedios.}

\begin{enumerate}

\item Área equivalente que soporta tensión axial.

\[\fitmath{y_{1}=2\cdot 140+2\cdot 80+3\cdot 140+3\cdot 80\approx 1100\,\mathrm{mm^{2}}}\]

\item Propiedades geométricas de la sección (centroide y momentos de inercia).

\[\fitmath{z_G=64{,}9090909\,\mathrm{mm},\quad y_G=45{,}0909091\,\mathrm{mm}}\]

\item Momentos de inercia del modelo en los ejes centroidales.

\[\fitmath{I_z=1304824{,}24\,\mathrm{mm}^{4},\quad I_y=3074824{,}24\,\mathrm{mm}^{4},\quad I_{yz}=28509{,}0909\,\mathrm{mm}^{4}}\]

\item Centro de cortadura a partir del equilibrio de los flujos; coordenadas del modelo.

\[\fitmath{z_C=59{,}1760692\,\mathrm{mm},\quad y_C=49{,}364584\,\mathrm{mm}}\]

\item Tensión máxima |q|/t en pared derecha.

\[\fitmath{y_{5}=\frac{\left|\frac{3996493470}{81946733}\right|}{2}\approx 24{,}3847\,\mathrm{MPa}}\]

\end{enumerate}
